\documentclass[11pt, a4paper]{article}
\usepackage[T1]{fontenc}
\usepackage[french]{babel}
\usepackage[margin=1in]{geometry}
\usepackage{graphicx}
\usepackage{amsmath}
\usepackage{amssymb}
\usepackage{hyperref}
\usepackage{xcolor}
\usepackage{listings}
\usepackage{fancyhdr}
\usepackage{booktabs}

\setlength{\headheight}{15pt}

\title{\textbf{Modele de Tracage des Connaissances Base sur les Graphes (GKT)} \\ \large Fonctionnement et Architecture Complete}
\author{Resume Detaille}
\date{\today}

\pagestyle{fancy}
\fancyhf{}
\rhead{GKT - Tracage des Connaissances par Graphes}
\lhead{Resume Technique}
\cfoot{\thepage}

\begin{document}

\maketitle

\tableofcontents
\newpage

\section{Introduction}

\subsection{Contexte et Motivation}

Le \textbf{tracage des connaissances} (knowledge tracing, KT) est un processus fondamental en apprentissage automatique applique a l'education. Il consiste a modeliser l'evolution de la maitrise d'un etudiant sur des concepts donnes en analysant ses interactions avec un systeme d'apprentissage intelligent (tuteur intelligent, plateforme en ligne, etc.).

Les modeles traditionnels comme :
\begin{itemize}
    \item \textbf{Bayesian Knowledge Tracing (BKT)} : modele probabiliste base sur les chaines de Markov cachees
    \item \textbf{Deep Knowledge Tracing (DKT)} : approche par reseaux de neurones recurrents (RNN/LSTM)
\end{itemize}

presentent des limitations importantes. Ils traitent essentiellement des \emph{sequences lineaires} d'interactions etudiantes, sans explicitement modeliser les \textbf{relations entre les concepts}.

\subsection{Innovation du Modele GKT}

Le modele \textbf{Graph-based Knowledge Tracing (GKT)} apporte une innovation majeure : il represente les relations entre concepts sous forme de \textbf{graphe de connaissances} et exploite des \textbf{reseaux de neurones a graphes} (Graph Neural Networks, GNN) pour propager l'information a travers cette structure.

\section{Principes Fondamentaux}

\subsection{Structure du Graphe de Connaissances}

Le coeur du modele GKT repose sur une representation graphique des concepts :

\begin{itemize}
    \item \textbf{Noeuds} : Chaque noeud represente un concept ou une competence (KnowledgeComponent, KC) en domaine d'apprentissage (p. ex. : addition, fractions, algebre, etc.)

    \item \textbf{Aretes} : Les aretes orientees entre noeuds representent les relations entre concepts, telles que :
          \begin{itemize}
              \item Relation de \emph{prerequis} (concept A doit etre maitrise avant concept B)
              \item Relation de \emph{co-requisit} (concepts lies qui doivent etre appris ensemble)
              \item Relation de \emph{similarite semantique} (concepts conceptuellement proches)
              \item Relation de \emph{transfert de connaissances} (apprendre A facilite l'apprentissage de B)
          \end{itemize}

    \item \textbf{Ponderations d'aretes} : Chaque arete peut avoir un poids representant la force de la relation
\end{itemize}

\textbf{Construction du graphe :}
\begin{enumerate}
    \item \emph{Approche symbolique} : Utilisation de mappings curriculaires manuels ou de domaines d'experts
    \item \emph{Approche data-driven} : Extraction automatique des relations a partir de matrices de similarite ou de correlations dans les donnees
    \item \emph{Approche hybride} : Combinaison des deux methodes
\end{enumerate}

\subsection{Etat Latent et Representation de la Maitrise}

Pour chaque concept (noeud), le modele maintient :

\begin{equation}
    \mathbf{h}_{i,t} \in \mathbb{R}^d
\end{equation}

ou :
\begin{itemize}
    \item $i$ est l'indice du concept
    \item $t$ est le pas de temps (apres la t-ieme interaction)
    \item $d$ est la dimension de l'embedding
    \item $\mathbf{h}_{i,t}$ est le \textbf{vecteur d'etat latent} representant la maitrise du concept i par l'etudiant au temps t
\end{itemize}

Cet etat latent encode implicitement la probabilite de maitrise du concept et evolue dynamiquement au cours des interactions.

\section{Architecture Complete du Modele GKT}

\subsection{Vue d'Ensemble du Pipeline}

Le fonctionnement du modele GKT suit ce pipeline :

\begin{enumerate}
    \item \textbf{Initialisation} : Initialiser les etats latents de tous les concepts
    \item \textbf{Traitement de l'interaction} : Pour chaque interaction etudiante :
          \begin{enumerate}
              \item Recuperer le concept $c_t$ concerne par la question
              \item Encoder la reponse de l'etudiant (correct/incorrect)
              \item Mettre a jour l'etat du concept concerne
              \item Propager l'information via le graphe a l'aide de GNN
          \end{enumerate}
    \item \textbf{Prediction} : Predire la probabilite de succes pour les questions futures
\end{enumerate}

\subsection{Etape 1 : Initialisation des Etats}

Au temps $t=0$, tous les concepts commencent par un etat initial :

\begin{equation}
    \mathbf{h}_{i,0} = \mathbf{w}_0 \in \mathbb{R}^d
\end{equation}

Cet etat peut etre :
\begin{itemize}
    \item Un vecteur aleatoire appris par le reseau
    \item Un vecteur uniforme (tous les concepts a etat neutre)
    \item Un vecteur initialise a partir de connaissances prealables
\end{itemize}

\subsection{Etape 2 : Mise a Jour par Interaction}

Supposons qu'a l'instant t, l'etudiant repond a une question portant sur le concept $c_t$ avec une reponse $a_t$ (correcte ou incorrecte).

\subsubsection{Encodage de l'Interaction}

L'interaction est encodee par :
\begin{equation}
    \mathbf{x}_t = \text{Embed}(c_t, a_t)
\end{equation}

ou Embed est une fonction d'embedding qui combine :
\begin{itemize}
    \item L'identifiant du concept (one-hot ou embedding)
    \item L'indicateur de reussite/echec
\end{itemize}

\subsubsection{Mise a Jour Locale du Concept}

L'etat du concept implique est d'abord mis a jour localement :

\begin{equation}
    \mathbf{h}'_{c_t,t} = \text{Update}(\mathbf{h}_{c_t,t-1}, \mathbf{x}_t)
\end{equation}

Le module d'update peut etre :
\begin{itemize}
    \item Une cellule GRU (Gated Recurrent Unit)
    \item Une cellule LSTM
    \item Un perceptron multicouche (MLP)
\end{itemize}

Exemple avec GRU :
\begin{equation}
    \mathbf{h}'_{c_t,t} = \text{GRU}(\mathbf{h}_{c_t,t-1}, \mathbf{x}_t)
\end{equation}

\subsubsection{Propagation via Reseaux de Neurones a Graphes}

L'etape cruciale du modele GKT est la \textbf{propagation de l'information a travers le graphe}. Les concepts lies au concept $c_t$ via des aretes dans le graphe sont affectes indirectement par l'apprentissage de $c_t$.

Cela simule l'effet de \textbf{transfert de connaissances} : maitriser un concept peut faciliter l'apprentissage de concepts lies.

\textbf{Utilisations courantes de GNN :}

\paragraph{Graph Convolutional Network (GCN)}

Pour chaque concept i et pas de temps t :

\begin{equation}
    \mathbf{h}_{i,t}^{(k+1)} = \sigma\left(\sum_{j \in \mathcal{N}(i)} w_{ij} \mathbf{W}^{(k)} \mathbf{h}_{j,t}^{(k)} + \mathbf{b}^{(k)}\right)
\end{equation}

ou :
\begin{itemize}
    \item $k$ est la couche GNN
    \item $\mathcal{N}(i)$ est l'ensemble des voisins du noeud i dans le graphe
    \item $w_{ij}$ est le poids normalise de l'arete (i,j)
    \item $\mathbf{W}^{(k)}$ est la matrice de poids apprenable
    \item $\sigma$ est une fonction d'activation (ReLU, etc.)
\end{itemize}

\paragraph{Graph Attention Network (GAT)}

Une variante qui apprend les poids d'attention :

\begin{equation}
    \alpha_{ij} = \text{softmax}_j\left(\text{LeakyReLU}(\mathbf{a}^T[\mathbf{W}\mathbf{h}_i \| \mathbf{W}\mathbf{h}_j])\right)
\end{equation}

\begin{equation}
    \mathbf{h}_{i,t}^{(k+1)} = \sigma\left(\sum_{j \in \mathcal{N}(i)} \alpha_{ij} \mathbf{W}^{(k)} \mathbf{h}_{j,t}^{(k)}\right)
\end{equation}

Les coefficients $\alpha_{ij}$ permettent au modele d'apprendre quels voisins sont les plus importants pour chaque concept.

\subsection{Etape 3 : Prediction}

Apres propagation, les etats latents mis a jour sont utilises pour predire la probabilite de succes :

\begin{equation}
    P(\text{succes}_{c,t+1} | \mathbf{h}_{c,t}) = \sigma(\mathbf{W}^T \mathbf{h}_{c,t} + b)
\end{equation}

ou avec un MLP plus complexe :

\begin{equation}
    P(\text{succes}_{c,t+1} | \mathbf{h}_{c,t}) = \sigma(\mathbf{W}_2 \text{ReLU}(\mathbf{W}_1 \mathbf{h}_{c,t} + \mathbf{b}_1) + \mathbf{b}_2)
\end{equation}

ou $\sigma$ est la fonction sigmoid assurant que la sortie est dans [0,1].

\section{Exemple Complet de Fonctionnement}

Considerons un exemple simplifie dans le domaine mathematique.

\subsection{Graphe de Concepts}

Supposons 4 concepts avec les relations suivantes :

\begin{center}
    \begin{tabular}{|c|c|c|}
        \hline
        \textbf{Concept} & \textbf{ID}  & \textbf{Description}      \\
        \hline
        0                & Arithmetique & Operations de base        \\
        1                & Fractions    & Manipulation de fractions \\
        2                & Algebre      & Resolution d'equations    \\
        3                & Geometrie    & Calcul de surfaces        \\
        \hline
    \end{tabular}
\end{center}

\textbf{Relations (Aretes)} :
\begin{itemize}
    \item Arithmetique $\rightarrow$ Fractions (prerequis)
    \item Arithmetique $\rightarrow$ Algebre (prerequis)
    \item Fractions $\rightarrow$ Algebre (complementaire)
    \item Algebre $\leftrightarrow$ Geometrie (concepts lies)
\end{itemize}

\subsection{Initialisation}

\begin{equation}
    \mathbf{h}_{0,0} = \mathbf{h}_{1,0} = \mathbf{h}_{2,0} = \mathbf{h}_{3,0} = [0, 0, 0, 0]^T
\end{equation}

Dimension d'embedding : $d = 4$

\subsection{Interaction 1 : L'Etudiant Repond a une Question d'Arithmetique}

L'etudiant repond \textbf{correctement} a une question sur l'Arithmetique.

\begin{enumerate}
    \item \textbf{Encodage} : $\mathbf{x}_1 = [1, 1, 0, 0]^T$ (concept 0, reponse correcte)

    \item \textbf{Mise a jour locale} :
          \begin{equation}
              \mathbf{h}'_{0,1} = \text{GRU}(\mathbf{h}_{0,0}, \mathbf{x}_1) \approx [0.3, 0.2, 0.1, 0.05]^T
          \end{equation}

          Les autres concepts restent inchanges initialement.

    \item \textbf{Propagation via GCN} :

          Le concept Arithmetique (0) affecte ses voisins Fractions (1) et Algebre (2) :

          \begin{align}
              \mathbf{h}_{1,1}^{\text{new}} & = \mathbf{h}_{0,0} + \alpha \cdot (\mathbf{W} \mathbf{h}'_{0,1}) \\
                                            & \approx [0.15, 0.1, 0.05, 0.02]^T
          \end{align}

          \begin{align}
              \mathbf{h}_{2,1}^{\text{new}} & = \mathbf{h}_{0,0} + \alpha \cdot (\mathbf{W} \mathbf{h}'_{0,1}) \\
                                            & \approx [0.15, 0.1, 0.05, 0.02]^T
          \end{align}

          Geometrie (3) n'est pas directement affectee.
\end{enumerate}

\subsection{Interaction 2 : L'Etudiant Repond a une Question de Fractions}

L'etudiant repond \textbf{incorrectement} a une question sur les Fractions.

\begin{enumerate}
    \item \textbf{Encodage} : $\mathbf{x}_2 = [0, 0, 1, 0]^T$ (concept 1, reponse incorrecte)

    \item \textbf{Mise a jour locale} :
          \begin{equation}
              \mathbf{h}'_{1,2} = \text{GRU}(\mathbf{h}_{1,1}, \mathbf{x}_2) \approx [0.1, 0.08, 0.04, 0.01]^T
          \end{equation}

          L'etat diminue car la reponse est incorrecte.

    \item \textbf{Propagation via GCN} :

          Fractions (1) affecte ses voisins Arithmetique (0) et Algebre (2).

          L'impact sur l'etat d'Algebre (2) est considere : bien que Fractions soit mal maitrisee, Algebre peut toujours progresser via Arithmetique.
\end{enumerate}

\section{Avantages du Modele GKT}

\begin{itemize}
    \item[$\checkmark$] \textbf{Capture les dependances} : Explicitement modelise les relations entre concepts

    \item[$\checkmark$] \textbf{Transfert de connaissances} : Simule comment maitriser un concept aide a en apprendre d'autres

    \item[$\checkmark$] \textbf{Interpretabilite} : La structure du graphe est explicite et comprehensible

    \item[$\checkmark$] \textbf{Generalisation} : Peut generaliser a des concepts non encore vus

    \item[$\checkmark$] \textbf{Performance} : Surpasse DKT et BKT sur plusieurs datasets

    \item[$\checkmark$] \textbf{Flexibilite} : Peut integrer differentes architectures GNN
\end{itemize}

\section{Limitations et Defis}

\begin{itemize}
    \item[$\times$] \textbf{Construction du graphe} : Requiert un effort pour creer un graphe precis

    \item[$\times$] \textbf{Scalabilite} : Peut devenir couteux computationnellement

    \item[$\times$] \textbf{Qualite des donnees} : Sensible a la qualite du graphe

    \item[$\times$] \textbf{Overfitting} : Avec peu de donnees, risque de surapprentissage

    \item[$\times$] \textbf{Parametres du graphe} : Determination des poids des aretes complexe
\end{itemize}

\section{Applications Pratiques}

\subsection{Tuteurs Intelligents Adaptatifs}

Les modeles GKT peuvent alimenter des systemes qui :
\begin{enumerate}
    \item Evaluent en temps reel la maitrise de chaque concept
    \item Recommandent les exercices optimaux bases sur les lacunes
    \item Expliquent les difficultES d'apprentissage
\end{enumerate}

\subsection{Analyse du Curriculum}

\begin{enumerate}
    \item Validation du graphe d'ordre
    \item Detection de gaps d'apprentissage
    \item Optimisation pedagogique
\end{enumerate}

\subsection{Evaluation Automatisee}

Predire le niveau de maitrise global d'un etudiant dans un domaine a partir d'interactions partielles.

\section{Extensions et Variantes}

\subsection{GKT Dynamique}

Adapter les poids des aretes pendant l'apprentissage :
\begin{equation}
    w_{ij}(t) = w_{ij}(t-1) + \lambda \cdot \Delta w_{ij}(t)
\end{equation}

Cela permet au modele de decouvrir les vraies relations entre concepts.

\subsection{GKT Multi-Modal}

Inclure des informations supplementaires :
\begin{itemize}
    \item Difficulte des questions
    \item Temps consacre aux taches
    \item Type d'exercice (QCM, reponse libre, etc.)
    \item Profil demographique de l'etudiant
\end{itemize}

\subsection{GKT Bidirectionnel}

Permettre aux aretes d'etre bidirectionnelles, modelisant les dependances mutuelles entre concepts.

\section{Comparaison avec les Modeles Traditionnels}

\begin{center}
    \begin{tabular}{|l|c|c|c|}
        \hline
        \textbf{Critere}   & \textbf{BKT} & \textbf{DKT} & \textbf{GKT}        \\
        \hline
        Modelise relations & Non          & Non          & \textbf{Oui}        \\
        Base sur graphe    & Non          & Non          & \textbf{Oui}        \\
        Interpretabilite   & Moyen        & Faible       & \textbf{Elevee}     \\
        Performance        & Bon          & Excellent    & \textbf{Excellent+} \\
        Complexite         & Basse        & Elevee       & Elevee              \\
        Donnees requises   & Moyen        & Eleve        & Moyen               \\
        \hline
    \end{tabular}
\end{center}

\section{Formulation Mathematique Complete}

\subsection{Notation et Cadre Probabiliste}

On considere une sequence d'interactions d'un etudiant :
\begin{equation}
    \mathcal{D} = \{(q_t, c_t, a_t)\}_{t=1}^{T}
\end{equation}
ou $q_t$ est la question, $c_t \in \{1,\dots,n\}$ le concept associe, et $a_t \in \{0,1\}$ la reponse (1 = correcte).

Le graphe de concepts est $G=(V,E)$ avec $|V|=n$ et matrice d'adjacence $\mathbf{A} \in \mathbb{R}^{n \times n}$. L'etat latent global est :
\begin{equation}
    \mathbf{H}_t = [\mathbf{h}_{1,t};\dots;\mathbf{h}_{n,t}] \in \mathbb{R}^{n \times d}
\end{equation}
avec $\mathbf{h}_{i,t}$ representant la maitrise du concept $i$ au temps $t$.

Le modele parametre une loi conditionnelle :
\begin{equation}
    P(a_t=1 \mid \mathcal{D}_{<t}, G; \Theta) = \hat{p}_t
\end{equation}
et l'apprentissage correspond a une estimation du maximum de vraisemblance reguliere.

\subsection{Normalisation du Graphe}

Pour stabiliser la propagation, on utilise classiquement :
\begin{equation}
    \tilde{\mathbf{A}} = \mathbf{A} + \mathbf{I}, \qquad
    \tilde{\mathbf{D}}_{ii} = \sum_j \tilde{\mathbf{A}}_{ij}
\end{equation}
et la matrice normalisee symetrique :
\begin{equation}
    \mathbf{S} = \tilde{\mathbf{D}}^{-1/2}\tilde{\mathbf{A}}\tilde{\mathbf{D}}^{-1/2}
\end{equation}
ce qui evite une explosion numerique sur les noeuds de fort degre.

\subsection{Mise a Jour Locale par GRU}

L'interaction est encodee par
\begin{equation}
    \mathbf{x}_t = \mathrm{Embed}(c_t) \oplus [a_t, 1-a_t].
\end{equation}

Pour le concept observe $c_t$, une cellule GRU realise :
\begin{align}
    \mathbf{z}_t         & = \sigma(\mathbf{W}_z\mathbf{x}_t + \mathbf{U}_z\mathbf{h}_{c_t,t-1} + \mathbf{b}_z),                     \\
    \mathbf{r}_t         & = \sigma(\mathbf{W}_r\mathbf{x}_t + \mathbf{U}_r\mathbf{h}_{c_t,t-1} + \mathbf{b}_r),                     \\
    \tilde{\mathbf{h}}_t & = \tanh(\mathbf{W}_h\mathbf{x}_t + \mathbf{U}_h(\mathbf{r}_t \odot \mathbf{h}_{c_t,t-1}) + \mathbf{b}_h), \\
    \mathbf{h}'_{c_t,t}  & = (1-\mathbf{z}_t) \odot \mathbf{h}_{c_t,t-1} + \mathbf{z}_t \odot \tilde{\mathbf{h}}_t.
\end{align}

Les autres noeuds sont conserves localement avant propagation :
\begin{equation}
    \mathbf{h}'_{i,t} = \mathbf{h}_{i,t-1}, \quad i \neq c_t.
\end{equation}

\subsection{Propagation Graphe sur L Couches}

En notant $\mathbf{H}'_t$ la matrice apres mise a jour locale, la propagation GCN s'ecrit :
\begin{align}
    \mathbf{H}_t^{(0)}      & = \mathbf{H}'_t,                                                                                                                      \\
    \mathbf{H}_t^{(\ell+1)} & = \phi\left(\mathbf{S}\mathbf{H}_t^{(\ell)}\mathbf{W}^{(\ell)} + \mathbf{1}(\mathbf{b}^{(\ell)})^\top\right), \quad \ell=0,\dots,L-1.
\end{align}

La sortie de propagation est $\mathbf{H}_t^{(L)}$. Cette operation couple explicitement les concepts voisins et formalise le transfert de connaissances.

\subsection{Tete de Prediction et Interpretation Probabiliste}

Pour predire la reponse future sur un concept cible $c$, on applique :
\begin{equation}
    \hat{p}_{c,t+1} = \sigma\left(\mathbf{w}_p^\top \mathbf{h}_{c,t}^{(L)} + b_p\right).
\end{equation}

Ainsi,
\begin{equation}
    a_{t+1} \mid \mathcal{D}_{\le t}, G \sim \mathrm{Bernoulli}(\hat{p}_{c,t+1}).
\end{equation}

\subsection{Objectif d'Optimisation}

Sur un ensemble de $M$ etudiants (sequence $m$ de longueur $T_m$), on minimise :
\begin{equation}
    \mathcal{L}(\Theta)
    = -\sum_{m=1}^{M}\sum_{t=1}^{T_m} \left[a_t^{(m)}\log \hat{p}_t^{(m)} + (1-a_t^{(m)})\log(1-\hat{p}_t^{(m)})\right]
    + \lambda \|\Theta\|_2^2.
\end{equation}

La derivee du terme BCE par rapport au logit $u_t$ (avec $\hat{p}_t=\sigma(u_t)$) est :
\begin{equation}
    \frac{\partial \mathcal{L}}{\partial u_t} = \hat{p}_t - a_t,
\end{equation}
ce qui facilite une optimisation stable par backpropagation a travers le temps et le graphe.

\subsection{Complexite Asymptotique}

Pour un pas de temps, le cout dominant est :
\begin{equation}
    \mathcal{O}(Ld|E| + Lnd^2 + d^2)
\end{equation}
ou $L$ est le nombre de couches GCN, $|E|$ le nombre d'aretes, et $d$ la dimension latente.

La memoire est de l'ordre :
\begin{equation}
    \mathcal{O}(nd + |E|)
\end{equation}
hors stockage des activations intermediaires pour l'apprentissage.

\section{Resultats Empiriques}

Sur plusieurs datasets publics (EdNet, KDD Cup 2010, etc.), les etudes montrent que GKT atteint :

\begin{itemize}
    \item \textbf{AUC} : 0.80-0.85 (vs 0.75-0.78 pour DKT)
    \item \textbf{Accuracy} : 78-82\% (vs 72-75\% pour BKT)
    \item \textbf{Temps d'inference} : Comparable ou legerement superieur a DKT
\end{itemize}

L'amelioration est particulierement notable pour les etudiants avec peu d'interactions (donnees clairsemees).

\section{Implementation Pratique}

\subsection{Frameworks Recommandes}

\begin{itemize}
    \item \textbf{PyTorch Geometric} : Excellente librarie pour GNN
    \item \textbf{TensorFlow GNN} : Alternative supportee par Google
    \item \textbf{PyTorch + DGL} : Deep Graph Library pour flexibilite maximale
\end{itemize}

\subsection{Etapes d'Implementation}

\begin{enumerate}
    \item Construire la matrice d'adjacence $\mathbf{A}$ du graphe
    \item Normaliser $\mathbf{A}$ (normalization D-A ou symetrique)
    \item Definir les couches GCN/GAT
    \item Creer la boucle d'entrainement avec backpropagation
    \item Evaluer sur un ensemble de test
\end{enumerate}

\section{Conclusion}

Le modele \textbf{Graph-based Knowledge Tracing (GKT)} represente une avancee significative dans la modelisation du processus d'apprentissage etudiant. En combinant la structure explicite de graphes de concepts avec la puissance des reseaux de neurones a graphes, GKT :

\begin{enumerate}
    \item Capture les \textbf{dependances complexes} entre concepts
    \item Fournit des \textbf{predictions plus precises} de maitrise
    \item Offre une \textbf{meilleure interpretabilite} pour les educateurs
    \item Simule le \textbf{transfert de connaissances} entre domaines connexes
\end{enumerate}

Bien que des defis demeurent (construction du graphe, scalabilite), GKT est prometteur pour la prochaine generation de \textbf{systemes d'apprentissage personnalises et adaptatifs}.

\begin{thebibliography}{99}

    \bibitem{wu2021} Wu, T., Chen, J., Zhang, Z., Niu, G., Zhang, Y., et Guestrin, C. (2021). \textit{Graph-based Knowledge Tracing: Modeling Student Learning with Concept Prerequisite Graph}. Proceedings of the AAAI Conference on Artificial Intelligence.

    \bibitem{piech2015} Piech, C., Bassen, J., Huang, J., Gagnauli, S., Sahami, M., Guibas, L. J., et Savarese, S. (2015). \textit{Deep knowledge tracing}. In Advances in Neural Information Processing Systems (pp. 505-513).

    \bibitem{corbett1994} Corbett, A. T., et Anderson, J. R. (1994). \textit{Knowledge tracing: Modeling the acquisition of procedural knowledge}. User modeling and user-adapted interaction, 4(4), 253-278.

    \bibitem{hamilton2017} Hamilton, W., Ying, Z., et Leskovec, J. (2017). \textit{Inductive representation learning on large graphs}. In Advances in Neural Information Processing Systems (pp. 1024-1034).

\end{thebibliography}

\end{document}